\documentclass[]{article}

\usepackage{amsmath}
\usepackage{multirow}
\usepackage{natbib}
\usepackage{graphicx} 


\begin{document}



In this paper, we tested the hypothesis that the formation of globular clusters (GCs) is regulated by the atomic cooling threshold, which posits that GCs assemble in dark matter halos with virial temperatures exceeding $10^4$ K. To test this framework across cosmic time, we utilized published data for two distinct GC populations: the 19 proto-GC candidates of the GEMS system observed at $z=9.625$ and the 5 GCs of the Sparkler system at $z=1.4$. For each of the 24 clusters, we calculated its formation redshift, $z_{form}$, from its published age and observed redshift. We then compared the distribution of these empirical formation epochs to the theoretical cosmic GC formation rate predicted by the atomic cooling model.

Our analysis yields two primary findings. First, the formation epochs of the two GC populations are fully consistent with a single, continuous theoretical formation history predicted by the atomic cooling model. The Sparkler GCs, with formation redshifts between $z_{form} \approx 2.2$ and $3.5$, align with the predicted peak of GC formation activity. In contrast, the GEMS GCs, with formation redshifts spanning $z_{form} \approx 9.7$ to $19.1$, populate the high-redshift tail of the same distribution. This is consistent with an observational selection effect, as only clusters that formed prior to the observation epoch of $z=9.625$ can be detected in the GEMS sample. The existence of these early clusters is in agreement with the non-zero formation rate predicted by the model at cosmic dawn.

Second, we find a striking and counter-intuitive result regarding the chemical enrichment of these systems. The GEMS clusters, despite forming at significantly earlier cosmic times, are systematically more metal-rich (median $[\mathrm{Z/H}] = -0.160$) than their lower-redshift Sparkler counterparts (median $[\mathrm{Z/H}] = -0.483$). This implies that the GEMS clusters were born within a host environment that underwent extremely rapid and efficient chemical enrichment in the first few hundred million years of the Universe. Furthermore, we identify a weak but statistically significant positive mass-metallicity relation within the GEMS system, suggesting that its more massive clusters are modestly more enriched.

From these results, we conclude that the atomic cooling threshold appears to be a universal and primary physical mechanism governing the formation of GCs from cosmic dawn to the peak of cosmic star formation. The successful alignment of two disparate GC populations, separated by over eight billion years of cosmic evolution, with a single theoretical framework provides strong evidence for this model. The high metallicities of the GEMS clusters offer a powerful constraint on their formation environment, pointing to their assembly within a massive and rapidly evolving host halo at $z>9.6$. This study demonstrates the unique power of high-redshift GC populations to probe the fundamental physics of galaxy assembly in the early Universe.

\end{document}
                